% Generated from data-plan.md; edit the Markdown, then rebuild.
% Markdown SHA256: a73a5ddba37a3f6a536275db56397825488b291714c69d8c8f1767c8b44e6569
% Shared template for Pandoc -> XeLaTeX/Tectonic. Content lives in the .md files.
\documentclass[11pt,a4paper]{article}
\usepackage[margin=22mm,headheight=14pt]{geometry}
\usepackage{fontspec}
\setmainfont{texgyretermes}[Extension=.otf,UprightFont=*-regular,BoldFont=*-bold,ItalicFont=*-italic,BoldItalicFont=*-bolditalic]
\setsansfont{texgyreheros}[Extension=.otf,UprightFont=*-regular,BoldFont=*-bold,ItalicFont=*-italic,BoldItalicFont=*-bolditalic]
\setmonofont{texgyrecursor}[Extension=.otf,UprightFont=*-regular,BoldFont=*-bold,ItalicFont=*-italic,BoldItalicFont=*-bolditalic,Scale=MatchLowercase]
\usepackage{amsmath,amssymb}
\usepackage{array,booktabs,longtable,calc}
\usepackage{graphicx}
\usepackage{xurl}
\usepackage[unicode,hidelinks]{hyperref}
\usepackage{bookmark}
\usepackage{enumitem}
\usepackage{titlesec}
\usepackage{fancyhdr}
\usepackage{microtype}
\hypersetup{pdftitle={Kế hoạch dữ liệu tuần
1},pdfauthor={Cao Văn Hà; Hà Mạnh Thế; Trần Duy Hưng},pdflang={vi-VN}}
\urlstyle{tt}
\setlength{\parindent}{0pt}
\setlength{\parskip}{5pt plus 1pt minus 1pt}
\linespread{1.08}
\setlist{nosep,leftmargin=*,topsep=4pt}
\setcounter{secnumdepth}{0}
\titleformat{\section}{\large\bfseries}{\thesection}{0pt}{}
\titleformat{\subsection}{\normalsize\bfseries}{\thesubsection}{0pt}{}
\titlespacing*{\section}{0pt}{12pt}{5pt}
\titlespacing*{\subsection}{0pt}{9pt}{4pt}
\setlength{\tabcolsep}{4pt}
\renewcommand{\arraystretch}{1.16}
\setlength{\LTpre}{6pt}
\setlength{\LTpost}{6pt}
\setlength{\emergencystretch}{2em}
\providecommand{\tightlist}{\setlength{\itemsep}{0pt}\setlength{\parskip}{0pt}}
\providecommand{\pandocbounded}[1]{#1}
\pagestyle{fancy}
\fancyhf{}
\fancyhead[L]{\small Đề tài 9 — Hồ sơ tuần 1}
\fancyhead[R]{\small Bản nháp}
\fancyfoot[C]{\thepage}
\renewcommand{\headrulewidth}{0.3pt}
\begin{document}
{\LARGE\bfseries Kế hoạch dữ liệu tuần 1\par}
\vspace{4pt}
{\large Kiến trúc Transformer và tinh chỉnh mô hình ngôn ngữ\par}
\vspace{4pt}
\textbf{Ứng dụng:} Hỏi đáp thủ tục hộ tịch–chứng thực tại Hà Nội.\\
\textbf{Nhóm:} Cao Văn Hà (trưởng nhóm), Hà Mạnh Thế, Trần Duy Hưng.\\
\textbf{Phiên bản:} Bản nháp tuần 1 — chưa có dữ liệu/huấn luyện đã kiểm chứng.
\vspace{5pt}\hrule\vspace{8pt}
\begin{quote}
\textbf{Trạng thái:} kế hoạch tuần 1, chưa phải datasheet của một bộ dữ
liệu đã thu thập. Các ngưỡng, mẫu thăm dò và quy trình dưới đây chuyển
yêu cầu của Đề tài 9 thành công việc có thể kiểm tra; chúng không chứng
minh dữ liệu, quyền dùng nguồn hay người tham gia đã sẵn sàng.
\end{quote}

\section{1. Mục đích, ranh giới và quyết định tiền điều
kiện}\label{mux1ee5c-ux111uxedch-ranh-giux1edbi-vuxe0-quyux1ebft-ux111ux1ecbnh-tiux1ec1n-ux111iux1ec1u-kiux1ec7n}

Bộ dữ liệu phục vụ ba mục tiêu tách biệt:

\begin{enumerate}
\def\labelenumi{\arabic{enumi}.}
\tightlist
\item
  huấn luyện/kiểm chứng mô hình ngôn ngữ ký tự nhỏ ở nhánh Phần A từ
  corpus có quyền dùng;
\item
  SFT trợ lý hỏi đáp thủ tục bằng câu hỏi--đáp có citation tới đoạn
  nguồn;
\item
  DPO từ cặp câu trả lời SFT cho cùng prompt, do tối thiểu hai người
  chấm độc lập đánh giá; và test tách biệt bằng câu hỏi người dùng thật
  nếu có consent.
\end{enumerate}

Đề tài không dùng PPO hay RAG. Không coi mô hình là nguồn pháp lý, không
thu hồ sơ cá nhân, không thu thông tin định danh để ``cá nhân hóa'' câu
trả lời. Dataset chỉ nhằm kiểm chứng nghiên cứu trong phạm vi thủ tục đã
công bố, không thay thế cổng/cơ quan có thẩm quyền.

Trước bất kỳ thu thập tự động hoặc khảo sát nào, nhóm phải kiểm tra ba
điều kiện vận hành:

\begin{enumerate}
\def\labelenumi{\arabic{enumi}.}
\tightlist
\item
  \textbf{Danh mục thủ tục và phiên bản nguồn nào} nằm trong phạm vi Hà
  Nội/hộ tịch--chứng thực đã thống nhất?
\item
  \textbf{Có thể tiếp cận người khảo sát với đồng ý bằng văn bản hay
  không?} Kế hoạch khảo sát đã được chấp thuận, nhưng chỉ gọi là test
  người dùng thật khi thực sự thu thập hợp lệ.
\item
  \textbf{Lịch chấm độc lập và phân xử thế nào?} Phương án Thế/Hưng chấm
  độc lập, Hà phân xử đã được đồng ý; cần hướng dẫn chấm và bố trí thời
  gian.
\end{enumerate}

Người dùng đã đồng ý phương án Hà Nội/hộ tịch--chứng thực, Truyện Kiều
cho mô hình ký tự, khảo sát sinh viên/người thân và bố trí Hà Mạnh
Thế--Trần Duy Hưng chấm độc lập, Cao Văn Hà phân xử. Đây là chấp thuận
kế hoạch, không phải bằng chứng đã có người tham gia, nhãn hay quyền thu
thập. Nhóm gồm Cao Văn Hà (trưởng nhóm, tổng hợp và nộp bài), Hà Mạnh
Thế, Trần Duy Hưng. Tài khoản/Drive, ngân sách, ngày nộp và mô hình nền
cuối vẫn mở; cần cụ thể hóa lịch chấm và quy trình đồng ý trước khảo
sát.

\subsection{1.1. Phương án dữ liệu đã được đồng
ý}\label{phux1b0ux1a1ng-uxe1n-dux1eef-liux1ec7u-ux111uxe3-ux111ux1b0ux1ee3c-ux111ux1ed3ng-uxfd}

\textbf{Đã được người dùng đồng ý về phương án:} địa bàn Hà Nội; đối
tượng trọng tâm là sinh viên/người trẻ và gia đình; nhóm thủ tục hộ
tịch--chứng thực. Bắt đầu mẫu nhỏ từ cấp bản sao trích lục hộ tịch/bản
sao giấy khai sinh và chứng thực bản sao từ bản chính; xem xét thêm
chứng thực chữ ký, xác nhận thông tin hộ tịch, cải chính/bổ sung hộ tịch
sau khi kiểm tra từng nguồn. Không coi các trường hợp cấp bản sao cho
từng loại giấy tờ là nhiều thủ tục độc lập chỉ để tăng số lượng. Chưa
đưa đất đai, tranh chấp, yếu tố nước ngoài vào phạm vi chính.

Nguồn đã khảo sát sơ bộ:

\begin{itemize}
\tightlist
\item
  \href{https://vpcp.dichvucong.gov.vn/p/home/dvc-tthc-thu-tuc-hanh-chinh-chi-tiet.html?ma_thu_tuc=1943}{Cổng
  Dịch vụ công quốc gia --- mã 2.000635}: đã đọc nội dung về hồ sơ, cách
  nộp, thời hạn, lệ phí và căn cứ pháp lý. Có trường trống và thông tin
  cấp thực hiện/địa chỉ cần đối chiếu; chưa coi toàn trang là đáp án
  chuẩn áp dụng tại Hà Nội.
\item
  \href{https://ochodua.hanoi.gov.vn/cai-cach-hanh-chinh/thu-tuc-chung-thuc-ban-sao-tu-ban-chinh-giay-to-van-ban-2832250922082409942.htm}{Cổng
  phường Ô Chợ Dừa --- chứng thực bản sao}: công cụ đọc hiện chỉ lấy
  được tiêu đề và mô tả, chưa xác minh đầy đủ nội dung thủ tục.
\item
  \href{https://ttpvhcc.hanoi.gov.vn/cai-cach-hanh-chinh-va-tin-khac/ha-noi-tai-cau-truc-thu-tuc-cap-ban-sao-trich-luc-ho-tich-ban-sao-giay-khai-sinh-tiet-kiem-hon-160-ty-dong-moi-nam-2850260102122029782.htm}{Trung
  tâm Phục vụ hành chính công Hà Nội --- thông tin tái cấu trúc thủ tục
  trích lục}: mới đọc metadata nhắc Quyết định 1818/QĐ-TTPVHCC; đây là
  đầu mối kiểm tra phiên bản địa phương, không thay văn bản quyết định
  gốc.
\end{itemize}

Đây là các đầu mối nghiên cứu, \textbf{không phải dữ liệu đã nhập kho
hay nguồn đã được duyệt quyền thu thập/phân phối}. Chưa kiểm
robots.txt/điều khoản từng cổng. Không lấy danh sách hồ sơ đã xử lý,
danh tính công dân hoặc thông tin nộp hồ sơ làm dữ liệu. Đối chiếu quy
định quốc gia với quy trình địa phương, ngày hiệu lực và hình thức nộp;
không chép lệ phí hay thời hạn giữa các địa phương.

Để đạt 2.000 QA có giá trị, đánh giá độ phủ sau mẫu 5--10 thủ tục/20--50
QA; nếu nguồn không đủ tình huống độc lập thì đề xuất mở rộng trong hộ
tịch--chứng thực, không nhân bản hàng trăm cách diễn đạt của cùng câu
hỏi. Chưa hứa số lượng thủ tục cần thiết trước kiểm kê.

\textbf{Corpus ký tự đã chọn về phương án:}
\href{https://vi.wikisource.org/wiki/Truy\%E1\%BB\%87n_Ki\%E1\%BB\%81u}{Truyện
Kiều trên Wikisource}. Đã đọc phần văn bản và nhãn phạm vi công cộng của
trang; liên kết phiên bản hiển thị \nolinkurl{oldid=206607}. Khi nhập
kho cần chốt revision, trích riêng thân tác phẩm, ghi công Nguyễn Du và
nguồn, kiểm điều kiện của bản số hóa, lưu hash; không lấy menu/chú thích
giao diện làm văn bản huấn luyện. Chia đoạn liên tục train/validation
trước khi tạo cửa sổ token; không để cửa sổ chồng lấp vượt ranh giới
split. Từ vựng học từ train, có cách xử lý ký tự chưa gặp. Văn bản thơ
nhỏ phù hợp mục tiêu kiểm chứng A nhưng dễ overfit và không đại diện văn
xuôi hành chính; PPL chỉ diễn giải trong corpus này. Không đưa văn bản
Kiều vào SFT hành chính.

\textbf{Khảo sát test đề xuất:} mời người từ 18 tuổi trở lên trong lớp,
câu lạc bộ và mạng lưới người thân có nhu cầu tìm hiểu thủ tục. Mục tiêu
vận hành ban đầu 50--80 người, mỗi người tự đặt 2--3 câu hỏi thật sự
muốn biết; tương ứng 100--240 câu thô trước lọc, không bảo đảm số test
hợp lệ. Đề không ấn định số người hay số câu test này. Không yêu cầu mỗi
người cố nghĩ đủ số câu, không đưa QA mẫu hoặc câu trả lời mô hình cho
họ chép. Thu tự nguyện, không gắn điểm học tập, không mặc định đồng ý vì
quen biết.

Biểu mẫu dự kiến hỏi: đồng ý tham gia sau thông tin mục đích/quyền rút;
đã đủ 18 tuổi hay chưa; từng làm/đang tìm hiểu nhóm thủ tục nào (không
hỏi số hồ sơ); câu hỏi muốn được giải đáp; câu hỏi có đang liên quan Hà
Nội hay không. Không thu tên, mã sinh viên, email tự động, địa chỉ cụ
thể hoặc ảnh giấy tờ. Xác nhận hình thức đồng ý điện tử đáp ứng yêu cầu
bằng văn bản trước thu; biểu mẫu không được lưu tự động nội dung nhạy
cảm chưa kiểm soát. Hướng dẫn người tham gia không nhập PII và có quy
trình ẩn danh trước lưu dữ liệu nghiên cứu. Cần quy định rõ thời hạn rút
trước khi phá liên kết định danh; sau ẩn danh hoàn toàn không hứa truy
tìm được dữ liệu từng người.

Khảo sát chủ yếu sinh viên là mẫu thuận tiện, không đại diện toàn bộ cư
dân Hà Nội. Giữ câu hỏi ngoài phạm vi phát sinh tự nhiên để đánh giá từ
chối; nếu không đủ, thu thêm câu hỏi thật cho nhóm này, không giả câu
hỏi AI là khảo sát. Khóa test và tách họ thủ tục theo thiết kế trước tạo
QA train; các câu bị trùng/rò rỉ phải xử lý trước huấn luyện, không sau
khi xem điểm.

\textbf{Phương án chấm DPO đã được đồng ý:} Hà Mạnh Thế và Trần Duy Hưng
chấm độc lập; Cao Văn Hà phân xử sau khi khóa nhãn và tính đồng thuận.
Ba người tập chấm một lô hiệu chuẩn riêng trước khi chấm chính thức;
lịch thực hiện cần xác định. Đề xuất sinh khoảng 650--750 cặp thô, giữ
ít nhất 500 cặp train hợp lệ và validation riêng sau loại hòa/lỗi; số
sinh này không bảo đảm đủ cặp hợp lệ. Đáp án phải sinh từ checkpoint SFT
thật, không viết tay chosen/rejected thay thế. Hai người cùng chấm mọi
cặp, không chia đôi 250 cặp/người. Chấm mù đánh giá cuối phải dùng tập
khác dữ liệu DPO; nếu không có người chấm ngoài nhóm, nêu giới hạn thiên
lệch nội bộ.

\section{2. Bốn nguồn/lớp dữ liệu và tiêu chí nhập
kho}\label{bux1ed1n-nguux1ed3nlux1edbp-dux1eef-liux1ec7u-vuxe0-tiuxeau-chuxed-nhux1eadp-kho}

{\def\LTcaptype{none} % do not increment counter
\begin{longtable}[]{@{}
  >{\raggedright\arraybackslash}p{(\linewidth - 8\tabcolsep) * \real{0.2000}}
  >{\raggedright\arraybackslash}p{(\linewidth - 8\tabcolsep) * \real{0.2000}}
  >{\raggedright\arraybackslash}p{(\linewidth - 8\tabcolsep) * \real{0.2000}}
  >{\raggedright\arraybackslash}p{(\linewidth - 8\tabcolsep) * \real{0.2000}}
  >{\raggedright\arraybackslash}p{(\linewidth - 8\tabcolsep) * \real{0.2000}}@{}}
\toprule\noalign{}
\begin{minipage}[b]{\linewidth}\raggedright
Mã lớp
\end{minipage} & \begin{minipage}[b]{\linewidth}\raggedright
Mục đích
\end{minipage} & \begin{minipage}[b]{\linewidth}\raggedright
Nguồn được phép và đơn vị thu thập
\end{minipage} & \begin{minipage}[b]{\linewidth}\raggedright
Điều kiện nhập kho
\end{minipage} & \begin{minipage}[b]{\linewidth}\raggedright
Không được suy diễn
\end{minipage} \\
\midrule\noalign{}
\endhead
\bottomrule\noalign{}
\endlastfoot
\nolinkurl{CHAR_CORPUS} & Huấn luyện LM ký tự nhỏ cho Phần A, báo cáo
PPL validation & Văn bản tiếng Việt thuộc phạm vi công cộng hoặc có giấy
phép mở, ví dụ Truyện Kiều \textbf{nếu} phiên bản/tình trạng quyền được
kiểm tra & Tài liệu có URL/định danh, giấy phép hoặc căn cứ phạm vi công
cộng, bản lưu/fingerprint và split train/val & Không dùng nội dung hành
chính chưa rõ quyền; không gọi corpus là dữ liệu Phần D \\
\nolinkurl{ADMIN_SOURCE} & Làm căn cứ cho QA/citation & Mô tả thủ tục
công khai tại \textbf{Cổng Dịch vụ công quốc gia và cổng của địa phương
đã chốt}, cùng văn bản/căn cứ công khai liên quan & Kiểm tra điều khoản,
robots.txt và giới hạn truy cập; lưu thời điểm, phiên bản/ngày hiệu lực
khi có, URL và span & Không bịa URL/cơ quan/văn bản; không lấy dữ liệu
cá nhân hay vượt kiểm soát truy cập \\
\nolinkurl{SFT_QA} và \nolinkurl{DPO_PAIR} & SFT, validation và DPO & QA
được tạo từ \nolinkurl{ADMIN_SOURCE}; hai ứng viên trả lời sinh từ
checkpoint SFT trên prompt đã khóa & Citation xác minh; SFT kiểm thủ
công theo mức tối thiểu từng split; preference có hai nhãn độc lập, rồi
mới phân xử & Không ghi nhận 2.000 QA hoặc 500 cặp khi chưa đủ count đã
kiểm; không dùng một người chấm thay cho hai người \\
\nolinkurl{USER_TEST} & Test độc lập và đánh giá mù & Câu hỏi do người
dùng thật cung cấp qua khảo sát sau consent bằng văn bản & Ẩn danh trước
lưu; tách hoàn toàn train/val; lưu consent-state không lưu định danh;
kiểm tra thuộc cụm người khảo sát & Không dùng câu hỏi ``tự nghĩ ra'' và
gọi là khảo sát thật; không giữ tên, liên hệ, hồ sơ hay dữ liệu nhạy cảm
không cần thiết \\
\end{longtable}
}

Cách triển khai ưu tiên là \textbf{đủ 2.000 QA train và 500 cặp DPO
train, cộng validation riêng}, vì train/validation không nên dùng chung
record. Đây là đề xuất bảo thủ để đáp ứng ngưỡng của đề bài sau chia
tách, không phải cam kết số lượng đã có. Quy mô \nolinkurl{USER_TEST} sẽ
chỉ được đặt sau khi biết khả năng khảo sát và số cụm người tham gia;
không làm đầy bằng dữ liệu mô phỏng.

Một pilot 5--10 thủ tục với 20--50 QA có thể giúp thử schema, citation,
split và rubric. Mẫu này \textbf{chưa thu thập}, phải kiểm danh mục
nguồn và quyền sử dụng trong phạm vi đã thống nhất trước khi thực hiện.
Pilot không thay thế bộ train yêu cầu và không phải ngưỡng nộp tuần 1
của giảng viên.

\section{3. Sổ nguồn, provenance và phiên
bản}\label{sux1ed5-nguux1ed3n-provenance-vuxe0-phiuxean-bux1ea3n}

Mọi item nguồn đi qua sổ nguồn trước khi sinh QA. Không scrape hàng loạt
trước khi kiểm điều kiện của từng domain và biểu mẫu nguồn.

\subsection{\texorpdfstring{3.1. Schema
\nolinkurl{source_document}}{3.1. Schema }}\label{schema-source_document}

{\def\LTcaptype{none} % do not increment counter
\begin{longtable}[]{@{}
  >{\raggedright\arraybackslash}p{(\linewidth - 4\tabcolsep) * \real{0.3333}}
  >{\raggedright\arraybackslash}p{(\linewidth - 4\tabcolsep) * \real{0.3333}}
  >{\raggedright\arraybackslash}p{(\linewidth - 4\tabcolsep) * \real{0.3333}}@{}}
\toprule\noalign{}
\begin{minipage}[b]{\linewidth}\raggedright
Trường
\end{minipage} & \begin{minipage}[b]{\linewidth}\raggedright
Kiểu/gợi ý
\end{minipage} & \begin{minipage}[b]{\linewidth}\raggedright
Ý nghĩa
\end{minipage} \\
\midrule\noalign{}
\endhead
\bottomrule\noalign{}
\endlastfoot
\nolinkurl{source_id} & UUID/chuỗi bất biến & Khóa nội bộ của một tài
liệu/phiên bản nguồn \\
\nolinkurl{source_type} & enum: \nolinkurl{char_corpus},
\nolinkurl{service_page}, \nolinkurl{legal_text}, \nolinkurl{guidance} &
Loại nguồn; không dùng tên mơ hồ \\
\nolinkurl{publisher_or_authority} & string & Tác giả/cơ quan phát hành
đúng như nguồn thể hiện \\
\nolinkurl{title} & string & Tiêu đề tài liệu/trang \\
\nolinkurl{canonical_url} & URL & URL chuẩn hoặc định danh công khai;
lưu nguyên văn \\
\nolinkurl{retrieved_at_utc} & ISO-8601 & Thời điểm nhóm truy xuất \\
\nolinkurl{published_at} / \nolinkurl{effective_from} /
\nolinkurl{effective_to} & nullable date & Ngày công bố/hiệu lực/hết
hiệu lực nếu nguồn công bố \\
\nolinkurl{source_version_label} & nullable string & Ký hiệu phiên
bản/số văn bản; không tự đặt thành ``mới nhất'' \\
\nolinkurl{content_hash} & SHA-256 & Fingerprint bản nội dung đã dùng \\
\nolinkurl{license_or_terms_url} & URL/text & Nơi kiểm tra giấy
phép/điều khoản dùng \\
\nolinkurl{robots_checked_at} & nullable ISO-8601 & Thời điểm kiểm
robots, nếu thu từ web \\
\nolinkurl{access_method} & enum + note & tải tay/API/crawl được phép;
rate-limit áp dụng \\
\nolinkurl{scope_group} & string & Nhóm thủ tục đã chốt để phục vụ
split \\
\nolinkurl{ingest_status} & enum & \nolinkurl{pending_review},
\nolinkurl{accepted}, \nolinkurl{rejected}, \nolinkurl{superseded} \\
\nolinkurl{reviewer_id} & pseudonymous ID & Người kiểm provenance/quyền;
không là PII công khai \\
\nolinkurl{notes} & text & Giới hạn, thay đổi hoặc lý do loại \\
\end{longtable}
}

Nếu phát hiện văn bản bị thay thế, tạo \nolinkurl{source_id}/phiên bản
mới hoặc liên kết \nolinkurl{supersedes_source_id}; không ghi đè bản cũ
rồi mất khả năng tái lập. Khi nguồn thiếu ngày hiệu lực, trường phải là
null và báo cáo nêu giới hạn.

\subsection{\texorpdfstring{3.2. Schema
\nolinkurl{source_span}}{3.2. Schema }}\label{schema-source_span}

Citation phải trỏ đến span bất biến của bản nguồn đã lưu, không chỉ đến
trang chủ.

{\def\LTcaptype{none} % do not increment counter
\begin{longtable}[]{@{}
  >{\raggedright\arraybackslash}p{(\linewidth - 4\tabcolsep) * \real{0.3333}}
  >{\raggedright\arraybackslash}p{(\linewidth - 4\tabcolsep) * \real{0.3333}}
  >{\raggedright\arraybackslash}p{(\linewidth - 4\tabcolsep) * \real{0.3333}}@{}}
\toprule\noalign{}
\begin{minipage}[b]{\linewidth}\raggedright
Trường
\end{minipage} & \begin{minipage}[b]{\linewidth}\raggedright
Kiểu/gợi ý
\end{minipage} & \begin{minipage}[b]{\linewidth}\raggedright
Ý nghĩa
\end{minipage} \\
\midrule\noalign{}
\endhead
\bottomrule\noalign{}
\endlastfoot
\nolinkurl{span_id} & UUID & Khóa span \\
\nolinkurl{source_id} & FK \nolinkurl{source_document} & Bản nguồn chứa
span \\
\nolinkurl{locator_type} & \nolinkurl{heading},
\nolinkurl{article_clause}, \nolinkurl{page}, \nolinkurl{char_offset},
\nolinkurl{xpath} & Cách xác định vị trí \\
\nolinkurl{locator_start}, \nolinkurl{locator_end} & string/int & Biên
span tái lập trong bản đã hash \\
\nolinkurl{heading_path} & nullable string & Đường dẫn đề
mục/Điều--Khoản khi có \\
\nolinkurl{quoted_text} & text & Đoạn trích đã chuẩn hóa tối thiểu để
review \\
\nolinkurl{normalization_version} & string & Quy tắc trích xuất/chuẩn
hóa áp dụng \\
\nolinkurl{span_hash} & SHA-256 & Fingerprint span \\
\nolinkurl{claim_tags} & array & Ví dụ: \nolinkurl{required_documents},
\nolinkurl{process}, \nolinkurl{time_limit}, \nolinkurl{fee},
\nolinkurl{legal_basis} \\
\nolinkurl{review_status} & enum & \nolinkurl{unreviewed},
\nolinkurl{verified}, \nolinkurl{rejected}, \nolinkurl{stale} \\
\end{longtable}
}

\nolinkurl{quoted_text} và bản nguồn có thể cần hạn chế quyền truy cập
nếu điều khoản nguồn yêu cầu; public manifest chỉ nên công khai URL,
locator và hash khi không được phép phân phối toàn văn.

\section{4. Schema cho QA SFT, DPO và
test}\label{schema-cho-qa-sft-dpo-vuxe0-test}

\subsection{\texorpdfstring{4.1. Schema
\nolinkurl{qa_record}}{4.1. Schema }}\label{schema-qa_record}

{\def\LTcaptype{none} % do not increment counter
\begin{longtable}[]{@{}
  >{\raggedright\arraybackslash}p{(\linewidth - 4\tabcolsep) * \real{0.3333}}
  >{\raggedright\arraybackslash}p{(\linewidth - 4\tabcolsep) * \real{0.3333}}
  >{\raggedright\arraybackslash}p{(\linewidth - 4\tabcolsep) * \real{0.3333}}@{}}
\toprule\noalign{}
\begin{minipage}[b]{\linewidth}\raggedright
Trường
\end{minipage} & \begin{minipage}[b]{\linewidth}\raggedright
Kiểu/gợi ý
\end{minipage} & \begin{minipage}[b]{\linewidth}\raggedright
Ý nghĩa
\end{minipage} \\
\midrule\noalign{}
\endhead
\bottomrule\noalign{}
\endlastfoot
\nolinkurl{qa_id} & UUID & Khóa bất biến \\
\nolinkurl{record_version} & semver/string & Phiên bản record sau chỉnh
sửa \\
\nolinkurl{split} & \nolinkurl{train}, \nolinkurl{validation},
\nolinkurl{test_user}, \nolinkurl{holdout} & Split khóa; test user không
được đưa lại train \\
\nolinkurl{scope_group} & string & Nhóm thủ tục để split theo nhóm \\
\nolinkurl{procedure_family_id} & string & Họ thủ tục; dùng phát hiện rò
rỉ cùng quy trình \\
\nolinkurl{question_text} & text & Prompt tiếng Việt đã loại PII \\
\nolinkurl{question_origin} & \nolinkurl{authoring},
\nolinkurl{user_survey}, \nolinkurl{paraphrase} & Nguồn gốc câu hỏi,
không nói mơ hồ ``thật'' \\
\nolinkurl{answer_text} & text & Đáp án SFT tham chiếu, có giới hạn/ từ
chối khi phù hợp \\
\nolinkurl{citation_span_ids} & array FK & Một hoặc nhiều span chống đỡ
từng mệnh đề chính \\
\nolinkurl{citation_rendering} & text & Dạng citation hiển thị cho mô
hình/người chấm \\
\nolinkurl{source_version_ids} & array FK & Phiên bản nguồn thực sự
dùng \\
\nolinkurl{claim_coverage} & structured array & Claim → span\_id → trạng
thái đủ/không đủ căn cứ \\
\nolinkurl{language} & \nolinkurl{vi} & Ngôn ngữ \\
\nolinkurl{provenance} & structured object & Tác giả/generator, thời
điểm, template, record nguồn \\
\nolinkurl{label_status} & enum & \nolinkurl{draft},
\nolinkurl{citation_checked}, \nolinkurl{approved},
\nolinkurl{rejected}, \nolinkurl{superseded} \\
\nolinkurl{labeler_id}, \nolinkurl{reviewer_id} & pseudonymous IDs &
Người tạo/kiểm; không đưa PII vào dataset \\
\nolinkurl{quality_flags} & array & \nolinkurl{needs_refusal},
\nolinkurl{stale_source}, \nolinkurl{ambiguous},
\nolinkurl{duplicate_candidate}\ldots{} \\
\nolinkurl{created_at}, \nolinkurl{updated_at} & ISO-8601 & Audit
trail \\
\end{longtable}
}

QA sinh bán tự động chỉ được chấp nhận sau review theo citation. Với
\textbf{mỗi split SFT có dùng cho huấn luyện/chọn cấu hình}, lấy mẫu
ngẫu nhiên phân tầng theo \nolinkurl{scope_group},
\nolinkurl{question_origin} và nguồn với khối lượng ít nhất
\(\lceil0{,}20N_{split}\rceil\); vì vậy riêng 2.000 QA train cần kiểm
thủ công ít nhất 400 record, và validation cũng có mẫu kiểm riêng theo
quy tắc này. Lưu \nolinkurl{qa_id}, reviewer, ngày, kiểm đủ
citation/đúng nội dung/đúng scope và quyết định. Nếu lỗi hệ thống xuất
hiện, mở rộng review hoặc sửa batch có liên quan; không chỉ thay vài mẫu
lỗi.

\subsection{\texorpdfstring{4.2. Schema
\nolinkurl{preference_pair}}{4.2. Schema }}\label{schema-preference_pair}

{\def\LTcaptype{none} % do not increment counter
\begin{longtable}[]{@{}
  >{\raggedright\arraybackslash}p{(\linewidth - 4\tabcolsep) * \real{0.3333}}
  >{\raggedright\arraybackslash}p{(\linewidth - 4\tabcolsep) * \real{0.3333}}
  >{\raggedright\arraybackslash}p{(\linewidth - 4\tabcolsep) * \real{0.3333}}@{}}
\toprule\noalign{}
\begin{minipage}[b]{\linewidth}\raggedright
Trường
\end{minipage} & \begin{minipage}[b]{\linewidth}\raggedright
Kiểu/gợi ý
\end{minipage} & \begin{minipage}[b]{\linewidth}\raggedright
Ý nghĩa
\end{minipage} \\
\midrule\noalign{}
\endhead
\bottomrule\noalign{}
\endlastfoot
\nolinkurl{pair_id} & UUID & Khóa cặp ưu tiên \\
\nolinkurl{pair_version} & string & Phiên bản cặp/rubric \\
\nolinkurl{split} & \nolinkurl{train}, \nolinkurl{validation},
\nolinkurl{test} & Tách validation DPO khỏi train; test chỉ khi thiết kế
cho phép \\
\nolinkurl{qa_id} / \nolinkurl{prompt_id} & FK & Cùng một prompt cho hai
ứng viên \\
\nolinkurl{prompt_text} & text/hash & Prompt chính xác đã sinh; có thể
tham chiếu nếu lặp \\
\nolinkurl{sft_checkpoint_id} & string/hash & Checkpoint sinh ứng
viên \\
\nolinkurl{generation_config} & JSON & temperature, top-p, max tokens,
seed và template \\
\nolinkurl{candidate_0_id}, \nolinkurl{candidate_1_id},
\nolinkurl{candidate_0_text}, \nolinkurl{candidate_1_text} & ID bất biến
+ text & Hai candidate theo thứ tự canonical, độc lập với A/B hiển thị
cho từng người chấm \\
\nolinkurl{candidate_0_citations}, \nolinkurl{candidate_1_citations} &
array FK/text & Citation người chấm có thể đối chiếu \\
\nolinkurl{rubric_version} & string & Phiên bản hướng dẫn chấm \\
\nolinkurl{rater_1_id}, \nolinkurl{rater_2_id} & pseudonymous IDs & Danh
tính vận hành tối thiểu của hai người chấm độc lập \\
\nolinkurl{rater_1_display_map}, \nolinkurl{rater_2_display_map} & JSON
& Ánh xạ A/B được hiển thị →
\nolinkurl{candidate_0_id}/\nolinkurl{candidate_1_id}; cho phép đảo thứ
tự riêng từng người \\
\nolinkurl{rater_1_label_raw}, \nolinkurl{rater_2_label_raw} &
\nolinkurl{A}, \nolinkurl{B}, \nolinkurl{tie}, \nolinkurl{invalid} &
Nhãn độc lập như người chấm đã thấy \\
\nolinkurl{rater_1_label_normalized},
\nolinkurl{rater_2_label_normalized} & \nolinkurl{candidate_0},
\nolinkurl{candidate_1}, \nolinkurl{tie}, \nolinkurl{invalid} & Nhãn quy
về candidate thực bằng display map; dùng cho agreement \\
\nolinkurl{rater_1_reason}, \nolinkurl{rater_2_reason} & coded/text & Lý
do theo rubric, hạn chế PII \\
\nolinkurl{rater_1_at}, \nolinkurl{rater_2_at} & ISO-8601 & Audit thời
điểm \\
\nolinkurl{independent_agreement} & derived & Đồng ý/bất đồng
\textbf{trước} phân xử, dựa trên nhãn normalized \\
\nolinkurl{adjudication_label} & \nolinkurl{candidate_0},
\nolinkurl{candidate_1}, \nolinkurl{exclude}, null & Quyết định sau bất
đồng \\
\nolinkurl{adjudicator_id}, \nolinkurl{adjudicated_at},
\nolinkurl{adjudication_reason} & audit fields & Ai/quy tắc đã phân
xử \\
\nolinkurl{chosen_text}, \nolinkurl{rejected_text} & derived & Chỉ tạo
sau khi có nhãn quyết định \\
\nolinkurl{provenance} & structured & Record/SFT checkpoint/nguồn liên
quan \\
\nolinkurl{status} & enum & \nolinkurl{pending}, \nolinkurl{labeled},
\nolinkurl{adjudicated}, \nolinkurl{excluded} \\
\end{longtable}
}

Hai người chấm phải xem candidate với thứ tự A/B đã đảo ngẫu nhiên, độc
lập và không biết nhãn của nhau. Họ dùng rubric: (i) đúng với span
nguồn, (ii) đầy đủ yêu cầu, (iii) citation đúng và có thể kiểm tra, (iv)
từ chối đúng khi ngoài phạm vi/thiếu căn cứ, (v) không bịa.
\nolinkurl{tie} và \nolinkurl{invalid} là kết quả hợp lệ; cặp như vậy
không được ép thành preference train. Trước phân xử, tính tỷ lệ đồng
thuận trên nhãn \nolinkurl{normalized} và Cohen's \(\kappa\) trên các
nhãn hợp lệ \nolinkurl{{candidate_0, candidate_1, tie}};
\nolinkurl{invalid} được báo riêng và loại khỏi mẫu tính \(\kappa\) theo
protocol này. Sau khi lưu nhãn độc lập, bất đồng mới được phân xử bằng
quy tắc đã công bố (ví dụ người thứ ba hoặc họp có biên bản); báo cáo tỷ
lệ đồng thuận, \(\kappa\), tỷ lệ loại và cách xử lý, không chỉ ghi nhãn
cuối.

\subsection{\texorpdfstring{4.3. Schema \nolinkurl{user_survey_record}
và tách dữ liệu con
người}{4.3. Schema  và tách dữ liệu con người}}\label{schema-user_survey_record-vuxe0-tuxe1ch-dux1eef-liux1ec7u-con-ngux1b0ux1eddi}

{\def\LTcaptype{none} % do not increment counter
\begin{longtable}[]{@{}
  >{\raggedright\arraybackslash}p{(\linewidth - 4\tabcolsep) * \real{0.3333}}
  >{\raggedright\arraybackslash}p{(\linewidth - 4\tabcolsep) * \real{0.3333}}
  >{\raggedright\arraybackslash}p{(\linewidth - 4\tabcolsep) * \real{0.3333}}@{}}
\toprule\noalign{}
\begin{minipage}[b]{\linewidth}\raggedright
Trường
\end{minipage} & \begin{minipage}[b]{\linewidth}\raggedright
Kiểu/gợi ý
\end{minipage} & \begin{minipage}[b]{\linewidth}\raggedright
Ý nghĩa
\end{minipage} \\
\midrule\noalign{}
\endhead
\bottomrule\noalign{}
\endlastfoot
\nolinkurl{survey_question_id} & UUID ngẫu nhiên & Khóa câu hỏi đã ẩn
danh \\
\nolinkurl{respondent_cluster_id} & UUID ngẫu nhiên, tách bảng liên hệ &
Cụm bootstrap, không dùng định danh thật \\
\nolinkurl{consent_version} & string & Phiên bản văn bản consent mà
người trả lời đã đồng ý \\
\nolinkurl{consent_status}, \nolinkurl{consented_at} & enum/time & Chỉ
nhận \nolinkurl{granted} trước khi ingest \\
\nolinkurl{withdrawn_at} & nullable time & Dấu vết xóa/loại nếu rút đồng
ý \\
\nolinkurl{question_text_redacted} & text & Câu hỏi đã xóa PII trước lưu
dataset \\
\nolinkurl{scope_group} & string & Kiểm tra thuộc phạm vi \\
\nolinkurl{survey_channel} & coarse enum & Kênh tổng quát, không lưu
liên hệ/địa chỉ chính xác \\
\nolinkurl{test_split_status} & \nolinkurl{eligible},
\nolinkurl{excluded}, \nolinkurl{heldout} & Khóa khỏi train/val/prompt
authoring \\
\nolinkurl{redaction_reviewer_id}, \nolinkurl{redaction_status} & audit
& Ai kiểm ẩn danh và trạng thái \\
\end{longtable}
}

Văn bản consent, thông tin liên hệ (nếu cần để vận hành khảo sát) và câu
hỏi đã ẩn danh phải lưu tách lớp truy cập. Dataset phân tích không chứa
tên, số điện thoại, địa chỉ, mã hồ sơ, ảnh giấy tờ, thông tin sức
khỏe/tài chính hoặc định danh tương tự. Nếu một câu hỏi chứa chi tiết
không thể ẩn danh mà vẫn giữ ý nghĩa, loại nó. Khi chưa có khảo sát hợp
lệ, \nolinkurl{USER_TEST} không được tạo bằng placeholder.

\section{5. Quy trình tạo dữ liệu và kiểm soát chất
lượng}\label{quy-truxecnh-tux1ea1o-dux1eef-liux1ec7u-vuxe0-kiux1ec3m-souxe1t-chux1ea5t-lux1b0ux1ee3ng}

\subsection{5.1. Ingest nguồn → span →
QA}\label{ingest-nguux1ed3n-span-qa}

\begin{enumerate}
\def\labelenumi{\arabic{enumi}.}
\tightlist
\item
  Chốt domain/cơ quan/nhóm thủ tục; kiểm điều khoản, robots.txt, tần
  suất và phạm vi truy cập trước ingest.
\item
  Lưu metadata \nolinkurl{source_document}, bản/manifest hoặc hash;
  không sửa im lặng. Phân đoạn theo đề mục, Điều--Khoản, trang hoặc
  offset có thể tái xác định.
\item
  Lập \nolinkurl{source_span}; người kiểm xác nhận span có thể trả lời
  claim nào.
\item
  Sinh/soạn QA theo template nhưng không vượt phạm vi span. Mỗi mệnh đề
  chính của đáp án phải liên kết span; nếu nguồn không đủ, đáp án phải
  giới hạn/từ chối thay vì suy luận.
\item
  Kiểm citation thủ công với khối lượng \(\lceil0{,}20N_{split}\rceil\)
  cho từng split SFT dùng cho train/validation; riêng 2.000 QA train tối
  thiểu là 400 record. Lưu kết quả review và sửa theo batch khi phát
  hiện lỗi chung.
\item
  Khóa split, tạo manifest số lượng/hashes; chỉ sau đó chuẩn bị dữ liệu
  huấn luyện.
\end{enumerate}

Nhóm không nên chờ đủ toàn bộ 2.000 QA mới thử pipeline. Cuối tuần 1
hoặc đầu tuần 2, sau khi có pilot đúng quyền và citation, chạy
\textbf{SFT pilot sớm} để phát hiện lỗi tokenizer/template/citation/VRAM
và để checkpoint SFT tạo ứng viên DPO. Pilot không được biến thành tuyên
bố rằng 500 cặp DPO đã tồn tại.

\subsection{5.2. Tạo preference cho
DPO}\label{tux1ea1o-preference-cho-dpo}

\begin{enumerate}
\def\labelenumi{\arabic{enumi}.}
\tightlist
\item
  Đóng băng checkpoint SFT pilot hoặc checkpoint đã chọn, tokenizer và
  prompt template.
\item
  Với cùng mỗi prompt, sinh hai candidate theo config đã ghi (có
  diversity đủ để nảy sinh lựa chọn); lưu seed/config/checkpoint, không
  sửa text sau khi thấy nhãn.
\item
  Randomize A/B và giao cho hai người đánh giá độc lập dùng rubric có ví
  dụ/calibration.
\item
  Khóa hai nhãn ban đầu, tính agreement; chỉ sau đó phân xử bất đồng
  theo quy tắc đã công bố. \nolinkurl{tie}, \nolinkurl{invalid} hoặc
  không đủ căn cứ bị loại khỏi train DPO và vẫn lưu lý do.
\item
  Tách train/validation DPO bằng khóa nhóm prompt/thủ tục, công bố count
  cuối theo \nolinkurl{status}; không trộn cặp cùng prompt qua split.
\end{enumerate}

Ngưỡng đề bài là ít nhất 500 cặp so sánh. Kế hoạch chỉ có thể coi ``500
train pairs'' đạt khi cặp đã qua hai nhãn, phân xử nếu cần, kiểm dữ liệu
và split. Tính \nolinkurl{pair_id} raw, số chấm độc lập, số đồng
thuận/bất đồng, số bị loại và số train/validation riêng; không báo một
con số gộp che mất cặp không dùng được.

\subsection{5.3. Kiểm soát trùng lặp, near duplicate và rò rỉ
split}\label{kiux1ec3m-souxe1t-truxf9ng-lux1eb7p-near-duplicate-vuxe0-ruxf2-rux1ec9-split}

\begin{itemize}
\tightlist
\item
  Chuẩn hóa Unicode (NFC), khoảng trắng và dạng citation trước khi so
  exact hash; lưu phiên bản chuẩn hóa.
\item
  Phát hiện near duplicate bằng chữ ký n-gram/MinHash hoặc embedding đã
  được chọn trong manifest, sau đó review thủ công các cặp qua ngưỡng.
  Không chỉ so text đáp án: phải xét prompt, procedure family, source
  span và paraphrase.
\item
  Split cấp \textbf{nhóm thủ tục/họ thủ tục} và khi phù hợp theo
  nguồn/phiên bản, thay vì random từng QA. \nolinkurl{holdout} là một
  hay nhiều họ thủ tục không có QA, source span hay paraphrase liên quan
  trong train/validation; nó đo khả năng khái quát sang tri thức thủ tục
  chưa thấy, không phải bằng chứng mô hình biết pháp lý mới.
\item
  Giữ toàn bộ \nolinkurl{USER_TEST} ngoài authoring, SFT, DPO train/val,
  tuning prompt và selection hyperparameter. Nếu người khảo sát gửi
  nhiều câu hỏi, tất cả phải nằm trong một
  \nolinkurl{respondent_cluster_id} và không được phân tán qua split.
\item
  Sau split chạy audit cross-split: exact hash, near-duplicate
  candidates, shared \nolinkurl{source_span_id}, shared
  \nolinkurl{procedure_family_id}, prompt/candidate DPO và
  \nolinkurl{respondent_cluster_id}. Chiến lược strict split này không
  được nới sau khi xem kết quả mô hình; nếu cần sửa xung đột trước
  train, tạo split/dataset version mới và ghi lý do.
\end{itemize}

\section{6. Split, phiên bản và quyền truy
cập}\label{split-phiuxean-bux1ea3n-vuxe0-quyux1ec1n-truy-cux1eadp}

Mỗi release dữ liệu phải có \nolinkurl{dataset_version}, ngày tạo, danh
sách file + SHA-256, phiên bản nguồn, quy tắc chuẩn hóa, nhãn và rubric;
manifest chia tập, số lượng theo trạng thái và changelog. Bất kỳ thay
đổi nào ảnh hưởng text, citation, nhãn, split hoặc nguồn tạo release
mới. Huấn luyện phải tham chiếu một \nolinkurl{dataset_version} bất
biến.

Quyền truy cập tối thiểu:

\begin{itemize}
\tightlist
\item
  \textbf{Raw/source có điều kiện:} người phụ trách dữ liệu; không công
  bố lại nếu điều khoản không cho phép.
\item
  \textbf{QA/DPO đã review:} nhóm thực hiện, có audit log thay đổi.
\item
  \textbf{Survey consent/liên hệ:} người được ủy quyền tối thiểu; không
  chuyển cho người chấm/huấn luyện.
\item
  \textbf{User test đã ẩn danh, khóa:} evaluator được chỉ định; không
  đưa cho người tune mô hình cho đến khi khóa thử nghiệm theo kế hoạch.
\end{itemize}

Nếu không thể phát hành raw data, repository vẫn phải chứa script thu
thập/biến đổi hợp lệ, schema, manifest/hash và hướng dẫn tái lập trong
phạm vi điều khoản. Datasheet cuối phải mô tả nguồn gốc, giấy phép, lọc,
PII/consent, split, giới hạn và cách yêu cầu xóa/rút dữ liệu con người.

\section{7. Kế hoạch đánh giá và thống
kê}\label{kux1ebf-houx1ea1ch-ux111uxe1nh-giuxe1-vuxe0-thux1ed1ng-kuxea}

\subsection{7.1. Đơn vị đánh
giá}\label{ux111ux1a1n-vux1ecb-ux111uxe1nh-giuxe1}

So sánh mô hình gốc, SFT và SFT+DPO trên cùng \nolinkurl{USER_TEST} khóa
nếu khảo sát hợp lệ. Người chấm mù model identity và thứ tự phương án;
họ kiểm tính đúng nội dung, đủ ý, citation đúng, ảo giác và từ chối
đúng. Vì hệ không có RAG và không cập nhật tri thức thời gian thực, điểm
tốt trên \nolinkurl{holdout} không cho phép suy ra câu trả lời đúng cho
văn bản pháp lý mới/chưa có trong nguồn; giới hạn này phải nêu trong báo
cáo. Nếu chưa có \nolinkurl{USER_TEST}, báo cáo trung thực rằng thiếu
tập test người dùng thật thay vì thay nó bằng train-derived QA.

Protocol khóa phải định nghĩa mẫu số trước khi chấm:

\begin{itemize}
\tightlist
\item
  \textbf{Độ chính xác mệnh đề thông tin (claim accuracy):} số mệnh đề
  được người chấm xác nhận đúng và được nguồn hỗ trợ / tổng mệnh đề
  thông tin trong các câu trả lời được chấm. Mệnh đề không xác định hoặc
  không đủ nguồn vẫn ở mẫu số, không ở tử số; báo riêng số lượng này.
  Câu trả lời không có mệnh đề thông tin được đánh giá ở chỉ số từ chối,
  không tự cho accuracy bằng 1.
\item
  \textbf{Tỷ lệ ảo giác theo câu trả lời:} số câu trả lời có ít nhất một
  mệnh đề thông tin không được nguồn hỗ trợ / tổng câu trả lời được
  chấm. Đây là định nghĩa vận hành ``không được nguồn hỗ trợ'', không
  khẳng định mọi mệnh đề đó sai trong thế giới thực; báo số câu bị loại
  do lỗi kỹ thuật riêng.
\item
  \textbf{Từ chối ngoài phạm vi:} số câu ngoài phạm vi hoặc không có căn
  cứ mà mô hình từ chối/điều hướng đúng / tổng câu thuộc nhóm đó.
  \textbf{Từ chối sai trong phạm vi:} số câu trong phạm vi, có căn cứ đủ
  và được protocol xem là trả lời được nhưng mô hình từ chối / tổng câu
  trong nhóm này; đây không được gộp thành ``từ chối đúng''.
\item
  \textbf{Win rate chấm mù theo cặp:} \((W+0{,}5T)/(W+T+L)\), trong đó
  \(W,T,L\) là số thắng, hòa, thua của một mô hình trên các cặp hợp lệ;
  báo số \nolinkurl{invalid}/loại riêng.
\end{itemize}

Mọi CI là \textbf{95\% bootstrap} theo cụm người khảo sát
(\nolinkurl{respondent_cluster_id}) khi một người có nhiều câu; nếu mỗi
câu thuộc người khác hoặc không có lặp người, đơn vị cluster chính là
câu hỏi. Nêu số câu, số cụm, số lần resample, statistic và seed
bootstrap; không coi các quan sát seed × prompt là độc lập.

\subsection{7.2. Năm run và chọn cấu
hình}\label{nux103m-run-vuxe0-chux1ecdn-cux1ea5u-huxecnh}

Đề bài yêu cầu trung bình ± độ lệch chuẩn trên ít nhất 5 run hoặc 5
fold. Đề xuất chạy 5 seed \textbf{độc lập} cho cấu hình cuối (cố định
dataset split và protocol); mỗi run phải có seed, checkpoint, config,
môi trường và log riêng. Với seed \(s\), chuỗi run là
\nolinkurl{SFT_s → DPO_s} với \nolinkurl{ref = SFT_s}; tập ưu tiên đã
khóa chung từ pilot, không phải chấm lại 500 cặp cho mỗi seed. Nhánh mô
hình ký tự Phần A cũng chạy ít nhất 5 seed độc lập để báo validation
PPL. Báo cáo mean và sample SD:

\[s=\sqrt{\frac{1}{n-1}\sum_{i=1}^{n}(x_i-\bar{x})^2},\qquad n=5.\]

Năm lần generation/decoding từ một checkpoint chỉ là biến thiên đầu ra,
\textbf{không} phải năm run. Baseline checkpoint/config cố định không
được gán sample SD bằng lặp output. Tương tự, không gộp các quan sát
seed × prompt như thể chúng độc lập để tạo CI hẹp; variation across
training run và uncertainty across questions là hai tầng khác nhau, cần
báo riêng hoặc dùng quy trình tổng hợp đã định trước.

Tất cả hyperparameter (mô hình nền, LoRA/QLoRA \(r\), \(\alpha\),
learning rate, epoch, length, DPO \(\beta\), generation template) chỉ
chọn bằng validation. Test khóa chỉ dùng một lần cho tổng kết. Nếu tính
toán không đủ để hoàn thành 5 run, nhóm phải xin điều chỉnh \textbf{với
giảng viên} và ghi rõ số run thực hiện cùng lý do; không ngầm giảm
chuẩn, không chỉ dựa vào đồng thuận nội bộ, và không thay 5 seed bằng 5
generations.

\subsection{7.3. Cầu nối A--D cần dữ liệu
thật}\label{cux1ea7u-nux1ed1i-ad-cux1ea7n-dux1eef-liux1ec7u-thux1eadt}

Để đáp ứng yêu cầu phần A tham gia so sánh hoặc dùng trực tiếp ở D, một
minibatch từ \nolinkurl{DPO_PAIR} Phần D đã review sẽ được token hóa
theo template SFT/DPO. Cả custom DPO và thư viện chuẩn nhận cùng
\nolinkurl{input_ids}, response token mask, candidate chosen/rejected,
policy/ref checkpoint và reduction. Với vanilla DPO,
\nolinkurl{sequence log-prob} là \textbf{tổng} log-prob token thuộc phần
trả lời đã mask và loss là trung bình theo ví dụ trong batch, không
chuẩn hóa trung bình token như PPL. Log artifact phải lưu
\nolinkurl{pair_id} (hoặc hash nếu dữ liệu hạn chế), length/mask, tổng
log-prob chosen/rejected, margin, loss, gradient đã chọn để đối chiếu và
sai khác số học.

Đây là đối chiếu trên log-prob và gradient từ forward/backward pass
\textbf{thật}, không dùng scalar điền tay. Với trạng thái tham số chung,
so sánh gradient autograd theo logit đã mask hoặc tham số adapter có thể
đối chiếu. Trước khi quy lỗi khác biệt, kiểm tokenizer, padding side,
mask prompt/response, truncation, label shift, beta, ref adapter và
batch reduction. Kết quả chỉ được kết luận sau run; kế hoạch này không
khẳng định sai số hay độ khớp đã đo.

\section{8. Rủi ro dữ liệu và điều kiện
dừng}\label{rux1ee7i-ro-dux1eef-liux1ec7u-vuxe0-ux111iux1ec1u-kiux1ec7n-dux1eebng}

{\def\LTcaptype{none} % do not increment counter
\begin{longtable}[]{@{}
  >{\raggedright\arraybackslash}p{(\linewidth - 4\tabcolsep) * \real{0.3333}}
  >{\raggedright\arraybackslash}p{(\linewidth - 4\tabcolsep) * \real{0.3333}}
  >{\raggedright\arraybackslash}p{(\linewidth - 4\tabcolsep) * \real{0.3333}}@{}}
\toprule\noalign{}
\begin{minipage}[b]{\linewidth}\raggedright
Rủi ro
\end{minipage} & \begin{minipage}[b]{\linewidth}\raggedright
Dấu hiệu
\end{minipage} & \begin{minipage}[b]{\linewidth}\raggedright
Hành động
\end{minipage} \\
\midrule\noalign{}
\endhead
\bottomrule\noalign{}
\endlastfoot
Nguồn thiếu quyền/robots hoặc thay đổi & Không xác minh được terms/phiên
bản; hash khác sau truy xuất & Không ingest/không phân phối; ghi
rejected/superseded và thay nguồn hợp lệ sau khi chốt phạm vi \\
Citation không hỗ trợ claim & Review phát hiện span không nói điều đáp
án khẳng định & Sửa/loại QA, gắn \nolinkurl{needs_refusal}; audit các
record chung template/span \\
QA/pair chưa đủ số lượng đã kiểm & Count
\nolinkurl{approved}/\nolinkurl{adjudicated} thấp & Pilot và điều phối
sớm; báo count thật; không đánh dấu hoàn thành ngưỡng \\
Thiếu hai người chấm & Không có nhãn độc lập & Dừng tạo train DPO từ cặp
đó; chốt người chấm trước khi mở rộng generation \\
Không có consent/khảo sát & Không có \nolinkurl{granted} hoặc PII chưa
redaction & Không lưu vào test; không tuyên bố test người dùng thật \\
Near duplicate/rò rỉ split & Audit phát hiện chung
prompt/span/family/cluster & Di chuyển/loại record, tạo split version
mới trước train; ghi quyết định \\
GPU/budget làm chậm pilot/5 run & OOM hoặc hết quota & Pilot SFT và DPO
ref riêng; log cấu hình; xin điều chỉnh minh bạch nếu không thể đáp ứng
5 run \\
\end{longtable}
}

\section{9. Deliverables tuần 1 và mốc tuần
2}\label{deliverables-tuux1ea7n-1-vuxe0-mux1ed1c-tuux1ea7n-2}

\textbf{Tuần 1:} chốt ba tiền điều kiện; tạo sổ nguồn/schema/rubric và
split policy; kiểm một mẫu nguồn hợp lệ; thử 5--10 thủ tục/20--50 QA chỉ
khi quyền và phạm vi đã được chốt; khởi động SFT pilot cuối tuần 1 nếu
có dữ liệu pilot đủ điều kiện.

\textbf{Cuối tuần 1--đầu tuần 2:} dùng checkpoint pilot để tạo candidate
DPO, chạy hai người chấm độc lập và thống kê agreement; tiếp tục tăng dữ
liệu SFT/DPO theo count đã review. Mốc tuần 2 của đề bài là bộ dữ liệu
Phần D cùng A/B cho báo cáo tiến độ; vì vậy cần ưu tiên thiết lập
pipeline preference sớm thay vì đợi đến sau khi đóng dữ liệu SFT.

Không có mục nào trong tài liệu này tuyên bố rằng 2.000 QA, 500 cặp,
khảo sát thật, checkpoint hay 5 run đã hoàn thành. Mọi số liệu cuối phải
dẫn về dataset/run manifest tương ứng.
\end{document}
